\scanpage{24}

Ici, \inserted{la sous-algèbre de Cartan est $\unsure{D|\mathfrak d|\delta}=\mathfrak t=kH$.}
\unsure{toute|tout} sous espace \deleted{de la forme} vectoriel de
$\mathfrak g$ contenant $\mathfrak t=kH$ \deleted{\unsure{la réponse du gpe à}} \inserted{est une sous-algèbre.}
\deleted{\unsure{faire varier le type}, $=$ Cartan}

Mais prenons $\mathfrak h=kH+k(\lambda X+\mu Y)$, si \deleted{$\mathfrak h$ était}
\deleted{\unsure{réalisé}} le normalisateur \inserted{$N$} de $\mathfrak h$ étant de dim 2,
\deleted{il contiendrait $T'$} \inserted{donc $\mathfrak h=\operatorname{Lie}(H)$ i.e. lisse} \deleted{donc} alors
\deleted{$H$ contiendrait un tore maximal $T'$, \unsure{donc}}
\deleted{\unsure{un conjugué de}} serait résoluble, donc un
\deleted{\unsure{visiblement facile}}\quad \deleted{le Borel}
sous-groupe de Borel, or \deleted{les seuls sous-groupes}
($T$ \unsure{étant|est} un C-sous-groupe) \deleted{\unsure{donc}} $\mathfrak h\supset\mathfrak t$

\deleted{le Borel-type $B$, tel que}
\deleted{$\mathscr L=kH+kX$, est tel que $B=\operatorname{Norm}\mathscr L$,}
\deleted{donc est \unsure{défini}}

impliquerait $H\supset T$,
\deleted{donc} or il n'y a que deux Borels
contenant $T$, \deleted{\unsure{si}} donc si $\lambda,\mu\ne0$, on
voit que $\unsure{\mathfrak h|\mathfrak g}$ n'est pas Algèbre de Lie d'un sous-groupe lisse de $G$ .....

3)

% archive sheet number: 23
